\chapter{Polars による式ベースの表処理}
\label{chap:polars}

Polars\index{polars@Polars|textbf} は、列指向の DataFrame ライブラリである。pandas と同じく表形式データを扱うが、index を持たず、列に対する処理を式として組み立てる点が異なる。\code{LazyFrame}\index{lazyframe@LazyFrame|textbf} では、式から作った実行計画を評価前に最適化できる。本章では、呼び出し時に結果を返す eager API と、処理を実行計画として蓄積する lazy API を区別し、列変換、集計、結合、入出力を扱う。

\section{式と遅延実行による表処理}

\subsection{index を持たない DataFrame}

Polars の DataFrame は、常に「行番号は整数位置でだけ参照する 2 次元表」である。pandas のような index に意味を持たせないため、\code{reset\_index()} や \code{loc} を前提にした設計にはならない。

\begin{lstlisting}[language=Python]
import polars as pl

df = pl.DataFrame(
    {
        "detector": ["A", "A", "B", "B"],
        "time_ns": [1032, 1188, 1055, 1244],
        "signal": [12.4, 10.1, 14.8, 9.7],
    }
)

print(df)
\end{lstlisting}

このコードを実行すると、Polars の表は次のように表示される。

\begin{lstlisting}[numbers=none, xleftmargin=2\zw]
shape: (4, 3)
+----------+---------+--------+
| detector | time_ns | signal |
| ---      | ---     | ---    |
| str      | i64     | f64    |
+==========+=========+========+
| A        | 1032    | 12.4   |
| A        | 1188    | 10.1   |
| B        | 1055    | 14.8   |
| B        | 1244    | 9.7    |
+----------+---------+--------+
\end{lstlisting}

\code{shape: (4, 3)} は 4 行 3 列であることを示し、各列名の下にはデータ型が表示される。Polars の式は行ラベルではなく列名を参照し、その列に適用する処理を表す。

\subsection{式による列変換}

Polars の中心概念は expression（式）\index{expression@expression|textbf}である。式は、列に対してどの計算を適用するかを表す処理記述であり、その時点の値そのものではない。\code{pl.col("signal")} は \code{signal} 列を参照する式であり、四則演算、比較、集計などと組み合わせられる。

\begin{lstlisting}[language=Python]
expr = pl.col("signal") / pl.col("signal").mean()
print(expr)
\end{lstlisting}

表示されるのは、概ね \code{[(col("signal")) / (col("signal").mean())]} のような式であり、4 行分の数値ではない。実際の値を得るには、\code{df.select(expr.alias("signal\_ratio"))} のように、式を DataFrame の処理へ渡す。\code{alias()} は式の結果へ列名を付ける。Polars の公式ドキュメント~\cite{polars_expressions} でも、expression は列処理を構成する基本概念として説明されている。\code{select()}、\code{with\_columns()}、\code{filter()}、\code{group\_by()} の多くは、式を受け取る API で統一されている。

\subsection{eager と lazy}
\index{polars@Polars!eager api@eager API}
\index{polars@Polars!lazy api@lazy API}

Polars には、命令ごとに結果を返す eager API と、処理を実行計画として蓄え、最後に最適化して実行する lazy API がある。

\begin{lstlisting}[language=Python]
eager_df = pl.read_csv("events.csv")
lazy_df = pl.scan_csv("events.csv")
\end{lstlisting}

\code{read\_csv()} は呼び出し時に読み込んで DataFrame を作るが、\code{scan\_csv()} は実行前の処理計画を表す \code{LazyFrame} を返す。lazy API では、必要な列だけを読む列射影や、不要な行を早い段階で除く述語プッシュダウンなどを適用できる場合がある。処理時間とメモリ使用量への効果は、入力形式と処理内容に依存するため、同じ入力で測定する。

\section{列変換・集計・結合・入出力}

\subsection{\code{select()} と \code{filter()}}

列の選択には \code{select()}\index{select@\code{select}|textbf}、条件に合う行の抽出には \code{filter()}\index{filter@\code{filter}|textbf} を用いる。

\begin{lstlisting}[language=Python]
selected = df.select("detector", "signal")
strong = df.filter(pl.col("signal") > 11.0)

print(selected)
print(strong)
\end{lstlisting}

\code{select()} で列が減り、\code{filter()} で条件に合う行だけが残る。出力は次の通りである。

\begin{lstlisting}[numbers=none, xleftmargin=2\zw]
shape: (4, 2)
+----------+--------+
| detector | signal |
+==========+========+
| A        | 12.4   |
| A        | 10.1   |
| B        | 14.8   |
| B        | 9.7    |
+----------+--------+
shape: (2, 3)
+----------+---------+--------+
| detector | time_ns | signal |
+==========+=========+========+
| A        | 1032    | 12.4   |
| B        | 1055    | 14.8   |
+----------+---------+--------+
\end{lstlisting}

\code{df["signal"]} は列を Series として取り出す。これに対して、\code{select()} や \code{filter()} へ渡した式は、複数列の変換や lazy API の実行計画にも使用できる。値を取り出す操作と、表に適用する式を区別する。

\subsection{\code{with\_columns()} による列の追加}

新しい列は \code{with\_columns()}\index{with_columns@\code{with\_columns}|textbf} で追加する。複数の式を一度に渡すと、それぞれの評価結果が列として追加または置換される。

\begin{lstlisting}[language=Python]
derived = df.with_columns(
    (pl.col("signal") / pl.col("signal").mean()).alias("signal_ratio"),
    (pl.col("time_ns") > 1100).alias("is_late"),
)

print(derived)
\end{lstlisting}

列追加後の表を見ると、式で作った列がどの位置へ入るかを確認できる。

\begin{lstlisting}[numbers=none, xleftmargin=2\zw]
shape: (4, 5)
+----------+---------+--------+--------------+---------+
| detector | time_ns | signal | signal_ratio | is_late |
+==========+=========+========+==============+=========+
| A        | 1032    | 12.4   | 1.055319     | false   |
| A        | 1188    | 10.1   | 0.859574     | true    |
| B        | 1055    | 14.8   | 1.259574     | false   |
| B        | 1244    | 9.7    | 0.825532     | true    |
+----------+---------+--------+--------------+---------+
\end{lstlisting}

\code{alias()} は式の結果へ列名を付ける操作である。式を複数組み合わせる場合は、出力列名を明示すると後続の選択や集計との対応を追いやすい。

\subsection{\code{group\_by()} と集計}

集計も式として書く。\code{group\_by()}\index{group_by@\code{group\_by}|textbf} と \code{agg()} は、pandas の \code{groupby()} に対応する中心的な道具である。

\begin{lstlisting}[language=Python]
summary = (
    derived
    .group_by("detector")
    .agg(
        pl.len().alias("count"),
        pl.col("signal").mean().alias("mean_signal"),
        pl.col("signal").std().alias("std_signal"),
        pl.col("is_late").mean().alias("late_fraction"),
    )
    .sort("detector")
)

print(summary)
\end{lstlisting}

集計後は、検出器ごとに 1 行ずつの表へまとまる。出力は次の通りである。

\begin{lstlisting}[numbers=none, xleftmargin=2\zw]
shape: (2, 5)
+----------+-------+-------------+------------+---------------+
| detector | count | mean_signal | std_signal | late_fraction |
+==========+=======+=============+============+===============+
| A        | 2     | 11.25       | 1.626346   | 0.5           |
| B        | 2     | 12.25       | 3.606245   | 0.5           |
+----------+-------+-------------+------------+---------------+
\end{lstlisting}

Polars の \code{group\_by()} は、集計対象が expression なので、同じ書き方のまま統計量を増やしやすい。\code{count}、\code{mean}、\code{std} のような基本操作は、式の名前を見れば意味を追いやすい。

\subsection{join と表の変形}
\index{polars@Polars!pivot@\code{pivot()}}

別表の較正係数や run 情報を結び付けるなら \code{join()}\index{join@\code{join}|textbf} を使う。縦長と横長の変換も行える。

\begin{lstlisting}[language=Python]
gain = pl.DataFrame(
    {
        "detector": ["A", "B"],
        "gain": [1.02, 0.97],
    }
)

merged = derived.join(gain, on="detector", how="left")
wide = (
    pl.DataFrame(
        {
            "detector": ["A", "A", "B", "B"],
            "kind": ["raw", "cal", "raw", "cal"],
            "value": [12.4, 12.7, 14.8, 14.4],
        }
    )
    .pivot(index="detector", on="kind", values="value")
    .sort("detector")
)

print(merged.select("detector", "signal", "gain"))
print(wide)
\end{lstlisting}

2 つの出力は次のようになる。

\begin{lstlisting}[numbers=none, xleftmargin=2\zw]
shape: (4, 3)
+----------+--------+------+
| detector | signal | gain |
+==========+========+======+
| A        | 12.4   | 1.02 |
| A        | 10.1   | 1.02 |
| B        | 14.8   | 0.97 |
| B        | 9.7    | 0.97 |
+----------+--------+------+
shape: (2, 3)
+----------+------+------+
| detector | raw  | cal  |
+==========+======+======+
| A        | 12.4 | 12.7 |
| B        | 14.8 | 14.4 |
+----------+------+------+
\end{lstlisting}

\code{join()} の \code{on="detector"} は両方の表で照合するキー列を指定し、\code{how="left"} は左側の \code{derived} の全行を残す結合方法を指定する。\code{pivot()} では、\code{index="detector"} が出力表で行を識別する列、\code{on="kind"} が新しい列名になる値、\code{values="value"} がセルへ入る値を表す。Polars には pandas のような index がないため、結合や変形に使う列を引数で明示する。

\subsection{欠損値、\code{null}、\code{NaN}}
\index{polars@Polars!null@\code{null}}
\index{nan@NaN|textbf}

Polars は欠損値の扱いが比較的厳密である。\code{null} と \code{NaN} を区別して考える。

\begin{lstlisting}[language=Python]
missing = pl.DataFrame(
    {
        "signal": [12.4, None, 14.8, float("nan")],
    }
)

print(missing.with_columns(
    pl.col("signal").is_null().alias("is_null"),
    pl.col("signal").is_nan().alias("is_nan"),
))
\end{lstlisting}

判定結果のうち、主要部分は次のようになる。

\begin{lstlisting}[numbers=none, xleftmargin=2\zw]
signal  is_null  is_nan
12.4    false    false
null    true     null
14.8    false    false
NaN     false    true
\end{lstlisting}

\code{null} は値が存在しないこと、\code{NaN} は浮動小数点の特殊値である。\code{is\_nan()} へ \code{null} を渡した結果も \code{null} になる点に注意する。両方を除くなら、例えば \code{pl.col("signal").is\_not\_null() \& pl.col("signal").is\_finite()} をフィルタ条件にする。どちらを落とすかによって使う式が変わる。

\subsection{文字列と条件分岐}

式ベースの API では、文字列処理や条件分岐も列全体に対する式として書く。

\begin{lstlisting}[language=Python]
names = pl.DataFrame(
    {
        "run_name": ["run_A_001", "run_A_002", "test_B_001"],
        "signal": [12.4, 10.1, 14.8],
    }
)

classified = names.with_columns(
    pl.col("run_name").str.extract(r"^(run_[A-Z])", 1).alias("prefix"),
    pl.when(pl.col("signal") > 12.0)
      .then(pl.lit("high"))
      .otherwise(pl.lit("normal"))
      .alias("quality"),
)

print(classified)
\end{lstlisting}

追加した 2 列を含む結果は次の通りである。

\begin{lstlisting}[numbers=none, xleftmargin=2\zw]
run_name    signal  prefix  quality
run_A_001   12.4    run_A   high
run_A_002   10.1    run_A   normal
test_B_001  14.8    null    high
\end{lstlisting}

\code{str.extract()} の第 1 引数は正規表現、第 2 引数の \code{1} は取り出すキャプチャグループの番号である。一致しない \code{test\_B\_001} の \code{prefix} は \code{null} になる。\code{when}、\code{then}、\code{otherwise} は、条件、真の場合の値、偽の場合の値を順に組み立てる。

\subsection{lazy API と \code{scan\_parquet()}}
\index{polars@Polars!scan_parquet@\code{scan\_parquet()}}
\index{polars@Polars!collect@\code{collect()}}

lazy API は、複数の読み込み、抽出、列選択、集計を一つの実行計画として最適化する場合に用いる。

\begin{lstlisting}[language=Python]
lf = (
    pl.scan_parquet("events/*.parquet")
    .filter(pl.col("signal") > 0.0)
    .select("detector", "signal")
    .group_by("detector")
    .agg(pl.col("signal").mean().alias("mean_signal"))
)

result = lf.collect()
print(result)
\end{lstlisting}

\code{scan\_parquet()} の引数 \code{"events/*.parquet"} は、一致する複数の Parquet ファイルを入力にする。戻り値は \code{LazyFrame} であり、この時点ではデータを実体化しない。\code{collect()} を呼ぶとクエリが実行され、戻り値として DataFrame が得られる。実行時には、\code{select()} で使う列だけを読む列射影や、\code{filter()} を早い段階へ移す最適化が適用されることがある。最適化内容を確認したい場合は、実行前に \code{print(lf.explain())} として計画を表示できる。

\subsection{CSV と Parquet の入出力}
\index{polars@Polars!parquet@Parquet}
\index{parquet@Parquet|textbf}

CSV の読み込み、Parquet への保存、遅延読み込みの型の違いを次の例で確認する。

\begin{lstlisting}[language=Python]
df = pl.read_csv("events.csv")
df.write_parquet("events.parquet")

lf = pl.scan_csv("events.csv")
result = lf.filter(pl.col("signal") > 0).collect()

print(type(df).__name__)
print(type(lf).__name__)
print(type(result).__name__)
\end{lstlisting}

表示される型名は順に \code{DataFrame}、\code{LazyFrame}、\code{DataFrame} である。\code{read\_csv()} はその場で読み込み、\code{scan\_csv()} は読み込み計画を作り、\code{collect()} が計画を実行する。\code{write\_parquet()} の第 1 引数は保存先であり、同名のファイルがある場合は上書きされるため、保存先を確認してから実行する。

\code{read\_csv()} と \code{write\_parquet()} は呼び出し時に入出力を行う。\code{scan\_csv()} は \code{LazyFrame} を返し、後続の式を実行計画へ追加する。計画の評価と DataFrame の生成は \code{collect()} が行う。

\section{pandas との違いと遅延実行の注意点}

\subsection{pandas の書き方をそのまま持ち込まない}

Polars と pandas では、同じ DataFrame という名称を使っていても API と実行モデルが異なる。Polars 公式の migration guide~\cite{polars_from_pandas} では、index を前提にせず、列式と処理文脈を用いる設計が説明されている。pandas の \code{loc} や \code{apply(axis=1)} を機械的に置き換えるのではなく、\code{select()}、\code{filter()}、組み込みの式で処理を表現する。

\subsection{\code{map\_elements()} や Python の \code{lambda} を多用しない}
\index{polars@Polars!map_elements@\code{map\_elements()}}

Polars の組み込み式は query optimizer が解析できる。Python の \code{lambda} や \code{map\_elements()} で任意の Python 関数を呼ぶと、最適化や並列実行の範囲が狭まり、呼び出しのオーバーヘッドも加わる場合がある。まず \code{pl.col()}、\code{when}、\code{str}、\code{dt}、\code{over} などの組み込み式で表せるかを確認し、表せない処理にだけ Python 関数を用いる。

\subsection{lazy API では \code{collect()} を忘れない}

\code{scan\_csv()} や \code{scan\_parquet()} は \code{LazyFrame} を返す。\code{print(lf)} が表示するのは実データではなく実行計画である。計画を評価して結果の DataFrame を得るには、\code{collect()} を呼び出す。

\subsection{欠損値における \code{null} と \code{NaN} の区別}

Polars では、各列のデータ型に応じて利用できる式が決まり、値が存在しない \code{null} と浮動小数点値 \code{NaN} は別に扱われる。フィルタや集計の前に列型を確認し、\code{is\_null()} と \code{is\_nan()} のどちらを条件に含めるかを決める。

\section{大規模表処理と可視化への受け渡し}

Polars への移行を検討する条件を次に示す。

\begin{itemize}
\item 複数の CSV または Parquet ファイルに対し、同じ列選択と行抽出を適用する場合
\item 列変換と集計を一つの実行計画として最適化する場合
\item pandas による入出力または集計が、測定した総所要時間の主要部分を占める場合
\item 列射影や述語プッシュダウンを適用できる入力形式と処理内容である場合
\end{itemize}

小規模データで index を用いた対話的操作が必要な場合や、既存コードが pandas の API に依存している場合は、pandas を維持する方が変更量を抑えられる。切り替えの判断には、行数だけでなく、必要な API、移行コスト、メモリ使用量、同じ入力に対する処理時間を用いる。

Polars の概念は Concepts~\cite{polars_concepts}、式は Expressions and contexts~\cite{polars_expressions}、遅延実行は Lazy API~\cite{polars_lazy}、入出力は IO~\cite{polars_io}、pandas との差は Coming from Pandas~\cite{polars_from_pandas} に記載されている。本章は、これらの機能を解析手順に沿って再構成した。

\section{検出器別集計と遅延実行}

共通入力には前章の \code{events.csv} を用いよ。ファイルがない場合は、第~\ref{chap:pandas}~章の章末問題に示した 6 行の CSV を作成せよ。

\begin{enumerate}
\item \textbf{列式による列選択と行抽出}。
\code{pl.read\_csv()} で表を読み、\code{select()} で \code{detector} と \code{signal} だけを選べ。次に、\code{signal} が \code{null} でなく 12.0 以上である行を \code{filter()} せよ。結果は A が 2 行、B が 1 行の計 3 行になる。\code{pl.col()} が値ではなく式を返すことと、\code{is\_not\_null()} を条件に明示する理由を説明せよ。

\item \textbf{列追加と集計を含む単一クエリ}。
\code{with\_columns()} で各検出器の平均からの差 \code{signal\_centered} を作成せよ。この計算には \code{pl.col("signal").mean().over("detector")} を用いよ。その後、\code{group\_by()} と \code{agg()} で件数、平均、最大値を求め、検出器名で並べ替えよ。\code{over("detector")} と \code{group\_by("detector")} の違いを、出力行数に着目して説明せよ。

\item \textbf{lazy API の実行計画と結果の確認}。
\code{pl.scan\_csv("events.csv")} から始め、必要な 2 列の選択、欠損除去、検出器別平均を記述せよ。\code{collect()} の前に \code{explain()} を表示し、その後に \code{collect()} した DataFrame を表示せよ。\code{LazyFrame} と \code{DataFrame} のどちらが実データを保持するかを説明せよ。

\item \textbf{結合と縦横変換}。
検出器と較正係数を持つ表を作り、\code{join(..., on="detector", how="left")} で結合して較正後信号を計算せよ。さらに、検出器を行、\code{valid} を列、件数を値とする表を \code{group\_by()} と \code{pivot()} で作れ。結合キーが重複している場合に出力行数が増える理由も説明せよ。

\item \textbf{pandas との結果照合}。
前章の検出器別集計と同じ統計量を Polars で求め、CSV に保存して pandas で読み戻せ。平均値が一致することを許容誤差付きで確認し、処理速度だけでなく、欠損値、列型、行順序も比較項目になる理由を述べよ。
\end{enumerate}

pandas と Polars で整形・集計した表は、数値を並べただけでは分布、外れ値、変数間の関係を把握しにくい。次章では、それらの列を Matplotlib へ渡し、解析途中の検証図と最終的な説明図を作成する。
